\chaptersection{2}{Identificazione degli asset}
DermaTriage comprende diversi asset che concorrono al processo di triage
e al mantenimento della pipeline nel tempo. Oltre ai dati clinici acquisiti
per ciascun caso, rientrano tra gli asset le informazioni prodotte durante
l'inferenza, i dataset e le strutture di persistenza, i modelli AI e gli
artefatti utilizzati nei processi di aggiornamento e validazione.

Le dipendenze tra questi elementi consentono agli effetti di un eventuale
malfunzionamento di propagarsi tra differenti stadi della pipeline. Gli asset
comprendono quindi sia le risorse direttamente coinvolte nella determinazione
della priorità clinica, sia quelle necessarie al recupero delle informazioni
e all'evoluzione del sistema.

\subsection{Dati clinici e input}
Il processo di triage riceve come input l'immagine della lesione e le informazioni cliniche associate al paziente, tra cui età, sesso, localizzazione anatomica e sintomatologia disponibile. Nel flusso integrato con B4, ogni richiesta è inoltre associata a un consulto specifico, il cui identificativo consente di mantenere il corretto collegamento tra i dati acquisiti, l'elaborazione effettuata e il risultato prodotto da DermaTriage.

% \begin{table}[!htbp]
%     \centering
%     \small
%     \setstretch{1}
%     \setlength{\tabcolsep}{4pt}
%     \renewcommand{\arraystretch}{1.05}

%     \begin{tabularx}{\textwidth}{
%         >{\raggedright\arraybackslash}p{0.6cm}
%         >{\raggedright\arraybackslash}p{2.5cm}
%         >{\raggedright\arraybackslash}p{1.5cm}
%         X
%         X
%         >{\raggedright\arraybackslash}p{1.8cm}
%     }
%         \toprule

%         \textbf{ID} &
%         \textbf{Asset} &
%         \textbf{Tipo} &
%         \textbf{Funzione} &
%         \textbf{Input / utilizzo} &
%         \textbf{Fase} \\

%         \midrule

%         A01 &
%         Immagine dermatologica &
%         Dato clinico &
%         Rappresenta visivamente la lesione cutanea sottoposta al processo
%         di triage. &
%         Acquisita tramite B4 o fornita direttamente tramite
%         \texttt{/analyze}; viene utilizzata da EfficientNet-B4 e
%         Qwen2-VL. &
%         Acquisizione e inferenza \\

%         \midrule

%         A02 &
%         Dati clinici, demografici e sintomatologia &
%         Dato clinico &
%         Fornisce il contesto clinico e demografico associato alla lesione
%         e completa le informazioni utilizzate durante l'elaborazione. &
%         Comprende età, sesso, localizzazione anatomica e sintomatologia
%         disponibile; i sintomi contribuiscono alla sintesi clinica
%         prodotta da BioMistral. &
%         Acquisizione e inferenza \\

%         \midrule

%         A03 &
%         Identificativi e associazione del consulto &
%         Metadato &
%         Mantiene il collegamento tra la richiesta di triage e i dati,
%         documenti e risultati appartenenti allo stesso caso. &
%         Nel percorso integrato con B4, il \texttt{consulto\_id} identifica
%         la richiesta da elaborare e permette di associare ad essa le
%         informazioni e il risultato diagnostico. &
%         Integrazione e gestione del caso \\

%         \bottomrule
%     \end{tabularx}

%     \caption{Asset relativi ai dati clinici e agli input di DermaTriage.}
%     \label{tab:asset-dati-input}
% \end{table}

% \begin{landscape}

% \begin{table}[!htbp]
%     \centering
%     \small
%     \setstretch{1}
%     \setlength{\tabcolsep}{4pt}
%     \renewcommand{\arraystretch}{1.08}

%     \begin{tabularx}{\linewidth}{
%         >{\raggedright\arraybackslash}p{0.7cm}
%         >{\raggedright\arraybackslash}p{3.0cm}
%         >{\raggedright\arraybackslash}p{2.2cm}
%         X
%         X
%         X
%         >{\raggedright\arraybackslash}p{2.4cm}
%     }
%         \toprule

%         \textbf{ID} &
%         \textbf{Asset} &
%         \textbf{Tipo} &
%         \textbf{Funzione} &
%         \textbf{Input / origine} &
%         \textbf{Output / destinazione} &
%         \textbf{Fase} \\

%         \midrule

%         A01 &
%         Immagine dermatologica &
%         Dato clinico grezzo &
%         Fornisce l'informazione visiva della lesione necessaria alla sua
%         classificazione e descrizione clinica. &
%         Immagine della lesione acquisita tramite B4 oppure fornita
%         direttamente tramite \texttt{/analyze}. &
%         Fornita a EfficientNet-B4 per la classificazione dell'urgenza e
%         a Qwen2-VL per la generazione della descrizione clinica. &
%         Acquisizione e inferenza \\

%         \midrule

%         A02 &
%         Dati clinici, demografici e sintomatologia &
%         Dato clinico strutturato &
%         Fornisce il contesto clinico e demografico complementare
%         all'immagine della lesione. &
%         Età, sesso, localizzazione anatomica e sintomatologia disponibile
%         associati al caso. &
%         Le informazioni contestuali vengono utilizzate durante
%         l'elaborazione; la sintomatologia costituisce un input della
%         sintesi clinica prodotta da BioMistral. &
%         Acquisizione e inferenza \\

%         \midrule

%         A03 &
%         Identificativo del consulto &
%         Metadato applicativo &
%         Identifica la richiesta B4 da elaborare e preserva l'associazione
%         tra il consulto e la relativa elaborazione. &
%         \texttt{consulto\_id} fornito a DermaTriage attraverso
%         \texttt{/diagnose}. &
%         Utilizzato per individuare il consulto B4 da elaborare e associare
%         il risultato prodotto alla richiesta corrispondente. &
%         Integrazione con B4 e gestione del caso \\

%         \bottomrule
%     \end{tabularx}

%     \caption{Asset relativi ai dati clinici e agli input di DermaTriage.}
%     \label{tab:asset-dati-input}
% \end{table}

% \end{landscape}

\begin{table}[!htbp]
    \centering
    \small
    \setstretch{1}
    \setlength{\tabcolsep}{5pt}
    \renewcommand{\arraystretch}{1.12}

    \begin{tabularx}{\textwidth}{
        >{\raggedright\arraybackslash}p{0.7cm}
        >{\raggedright\arraybackslash}p{3.0cm}
        >{\raggedright\arraybackslash}p{2.2cm}
        X
        >{\raggedright\arraybackslash}p{2.4cm}
    }
        \toprule

        \textbf{ID} &
        \textbf{Asset} &
        \textbf{Tipo} &
        \textbf{Funzione} &
        \textbf{Fase} \\

        \midrule

        A01 &
        Immagine dermatologica &
        Dato clinico grezzo &
        Fornisce l'informazione visiva della lesione necessaria alla sua
        classificazione e descrizione clinica. &
        Acquisizione e inferenza \\

        \addlinespace[2pt]

        \multicolumn{2}{@{}l}{\textbf{Input / origine}} &
        \multicolumn{3}{>{\raggedright\arraybackslash}X}{
            Immagine della lesione acquisita tramite B4 oppure fornita
            direttamente tramite \texttt{/analyze}.
        } \\

        \addlinespace[2pt]

        \multicolumn{2}{@{}l}{\textbf{Output / destinazione}} &
        \multicolumn{3}{>{\raggedright\arraybackslash}X}{
            Fornita a EfficientNet-B4 per la classificazione dell'urgenza
            e a Qwen2-VL per la generazione della descrizione clinica.
        } \\

        \midrule

        A02 &
        Dati clinici, demografici e sintomatologia &
        Dato clinico strutturato &
        Fornisce il contesto clinico e demografico complementare
        all'immagine della lesione. &
        Acquisizione e inferenza \\

        \addlinespace[2pt]

        \multicolumn{2}{@{}l}{\textbf{Input / origine}} &
        \multicolumn{3}{>{\raggedright\arraybackslash}X}{
            Età, sesso, localizzazione anatomica e sintomatologia
            disponibile associati al caso.
        } \\

        \addlinespace[2pt]

        \multicolumn{2}{@{}l}{\textbf{Output / destinazione}} &
        \multicolumn{3}{>{\raggedright\arraybackslash}X}{
            Le informazioni contestuali vengono utilizzate durante
            l'elaborazione; la sintomatologia costituisce un input della
            sintesi clinica prodotta da BioMistral.
        } \\

        \midrule

        A03 &
        Identificativo del consulto &
        Metadato applicativo &
        Identifica la richiesta B4 da elaborare e preserva l'associazione
        tra il consulto e la relativa elaborazione. &
        Integrazione con B4 e gestione del caso \\

        \addlinespace[2pt]

        \multicolumn{2}{@{}l}{\textbf{Input / origine}} &
        \multicolumn{3}{>{\raggedright\arraybackslash}X}{
            \texttt{consulto\_id} fornito a DermaTriage attraverso
            \texttt{/diagnose}.
        } \\

        \addlinespace[2pt]

        \multicolumn{2}{@{}l}{\textbf{Output / destinazione}} &
        \multicolumn{3}{>{\raggedright\arraybackslash}X}{
            Utilizzato per individuare il consulto B4 da elaborare e
            associare il risultato prodotto alla richiesta corrispondente.
        } \\

        \bottomrule
    \end{tabularx}

    \caption{Asset relativi ai dati clinici e agli input di DermaTriage.}
    \label{tab:asset-dati-input}
\end{table}